%% ===========================================================================
%% Rewrite of \subsubsection{Carry Handling in Microcode}
%%
%% Provenance of every claim below (see carry_handling_CLAUSE_MAP.md):
%%   master probe            Paper_probes/additionwithflags.c        (S1-S8)
%%   register-class scope    Paper_probes/carrystate_scope.c         (A,B,C,D)
%%   bridge operand matrix   Paper_probes/gfl_rr_operand_matrix.c
%%   bridge discovery        Paper_probes/genflagsrr_exhaustive.c
%%   false-positive killer   Paper_probes/test_genarithflags_semantics.c
%%   arch-dest failure       Paper_probes/arch_gfl_chain.c
%%   chain length            Paper_probes/chain_length.c, Paper_probes/chained_4limb_add.c
%%   intra-triad packing     Paper_probes/test_adc_gfl_adc_intratriad.c
%%   dead ends               Paper_probes/test_adc_chain.c, test_movetocreg_eflags.c,
%%                           cr_cf_probe.c
%%   arch CF frozen          Paper_probes/test_adc_dsz64.c, intra-triad-adc.c,
%%                           does_add_update_eflags.c, readaflags_probe.c
%% ===========================================================================

\subsubsection{Carry Handling in Microcode}
\label{sec:c25519-carries}

The implementation depends on how Goldmont microcode represents and consumes
arithmetic carry state.
This mechanism differs from the conventional x86 interface and motivates the
arithmetic organization used by both Curve25519 implementations.

An arithmetic $\mu$op such as \code{ADD} associates conditional state with its
destination \emph{register}, provided that destination is one of the sixteen
internal \code{TMP} registers.
For an addition, this state records whether the operation produced a carry;
it is in fact the whole arithmetic flag set, since \code{SETCC} variants
recover zero, sign and overflow from the same destination as well.
The state is per register and survives across triads: it persists until the
next $\mu$op writes that register, whatever that $\mu$op is --- a shift
replaces it with its own shifted-out bit, and a multiply replaces the state of
both registers it writes --- so scheduling must treat any write to a register
as destroying a carry still owed from it.
Several carries can nonetheless be live simultaneously in different registers
and be recovered in any order.
An architectural destination gets the correct arithmetic result but no
queryable condition state, which is the single most consequential constraint
on patch register allocation: every carry-producing addition in our patches
writes a \code{TMP}, and results are copied out to architectural registers
with \code{ZEROEXT} afterwards.

The patch can consume this state in two ways.

The first path uses \code{SETCC\_CONDB}.
This operation reads the conditional state associated with a previous
arithmetic destination and materializes the carry as an ordinary integer equal
to zero or one.
The restriction to internal registers applies only to the register whose state
is being read: \code{SETCC\_CONDB} may write its result straight to an
architectural register, as a full width zero extending write, and a subsequent
\code{ADD} accumulates that value into any register, architectural or
internal.
The carry therefore becomes an explicit register value rather than a flag,
which costs three $\mu$ops per propagated carry
(\code{ADD}, \code{SETCC\_CONDB}, folding \code{ADD})
but never touches \code{RFLAGS} and never serializes two carries against a
single shared flag register.
Different additions can materialize their carries into different registers, so
several accumulation chains can progress independently.

The second path reconstructs the conventional x86 carry mechanism.
A microcode \code{ADC} does not read the conditional state associated with an
arithmetic destination; it reads the architectural carry flag, which is
captured from \code{RFLAGS} when the hook fires and is then frozen for the
lifetime of the patch.
Neither \code{ADD} nor a preceding \code{ADC} updates it, so two \code{ADC}s
placed back to back both observe the same entry value.
To use a carry produced inside the patch, the patch must first publish it with
the two-operand form \code{GENARITHFLAGS\_RR(t,t)}, naming the \code{TMP} whose
carry is wanted in the first operand; the following \code{ADC} then adds its
two operands together with the architectural carry flag.
The operand form matters.
The single-operand \code{GENARITHFLAGS\_R} does not bridge: it writes a
different architectural flag image, one that \code{READAFLAGS} observes but
\code{ADC} ignores.\footnote{An earlier investigation of ours reported
\code{GENARITHFLAGS\_R} as the bridge on the strength of a 128-bit addition
that passed only inputs of the form $2^{64}-1+1$, where carry-out and a zero
low half coincide. Inputs that separate the two disprove it.}
The publishing register must also be a \code{TMP}: the same chain built on
architectural destinations silently drops carries.
Since \code{ADC}'s own destination latches a fresh condition, chains compose;
we verified chains of this shape to $32$ limbs, at roughly eight cycles per
limb amortized.
Each hop needs its own \code{GENARITHFLAGS\_RR}, so the path costs two $\mu$ops
per propagated carry, and the bridge and the next \code{ADC} can share a triad
with the addition that feeds them.
We found no other write path into the register \code{ADC} reads, having probed
\code{MOVEINSERTFLGS}, \code{MOVEMERGEFLGS} and all six \code{MOVETOCREG}
variants against \code{CORE\_CR\_EFLAGS}.

Table~\ref{tab:c25519-carrypaths} summarizes the two paths.

\begin{table}[htbp]
\centering
\caption{Two ways to consume carry state in Goldmont microcode.
The main $5\times51$ implementation materialises carries as integers and
accumulates them explicitly.
The saturated $4\times64$ control instead publishes carry state to the
architectural carry flag and consumes it with \code{ADC}.}
\label{tab:c25519-carrypaths}
\footnotesize
\setlength{\tabcolsep}{5pt}
\begin{tabularx}{\linewidth}{@{}>{\raggedright\arraybackslash}p{.22\linewidth}>{\raggedright\arraybackslash}X>{\raggedright\arraybackslash}X@{}}
\toprule
Operation path & Effect on carry state & Use in our implementations \\
\midrule
\code{ADD} $\rightarrow$ \code{SETCC\_CONDB} $\rightarrow$ \code{ADD}
\newline (3 $\mu$ops per carry)
&
An arithmetic $\mu$op records conditional state with its \code{TMP}
destination.
\code{SETCC\_CONDB} reads that destination associated state and materialises
the carry as an integer equal to zero or one.
A later \code{ADD} accumulates this value into any register.
The architectural flags are never touched, and carries in different registers
are independent.
&
Used in the main $5\times51$ multiplication and squaring patches
($25$ and $15$ occurrences respectively, no \code{ADC}).
This keeps the low sum, high sum, and carry count as separate dependency
chains. \\
\midrule
\code{ADD} $\rightarrow$ \code{GENARITHFLAGS\_RR} $\rightarrow$ \code{ADC}
\newline (2 $\mu$ops per carry)
&
\code{ADC} reads the architectural carry flag, which is frozen at hook entry
and updated by no arithmetic $\mu$op.
\code{GENARITHFLAGS\_RR(t,t)} publishes \code{TMP} $t$'s recorded carry into
it, and the next \code{ADC} consumes it.
Every destination in the chain must be a \code{TMP}; architectural
destinations drop the carry.
&
Used in the saturated $4\times64$ control.
This reconstructs a conventional full width carry chain, at the cost of
staging accumulators through \code{TMP} registers. \\
\bottomrule
\end{tabularx}
\end{table}

This distinction interacts with the choice of field representation.
The $5\times51$ representation leaves thirteen unused bits in each 64 bit limb,
so a column sum absorbs many products before any carry has to leave it: carries
are comparatively rare, and paying three $\mu$ops to move one as a value costs
little.
The saturated $4\times64$ representation instead propagates between adjacent
words on every column, which makes the per-carry cost the dominant term and
makes the two $\mu$op path worth its \code{TMP} staging overhead.
We built the $4\times64$ control both ways to confirm this rather than assume
it.
The \code{SETCC} form occupies $119$ triads and runs at $215$ cycles per
multiplication; the \code{ADC} form, once the bridge and the following
\code{ADC} were packed into shared triads, occupies $75$ triads and runs at
$200$ cycles.
The control we report uses the \code{ADC} form.
Both remain far behind the $58$ cycle $5\times51$ multiplication, so the
representation, not the carry path, is what decides the comparison.
